% ============================================================================
\chapter{The 16 Two-Input Boolean Functions}
\label{app:gates}
% ============================================================================

This appendix provides the complete reference for all 16 two-input
Boolean functions. Each function is identified by its 4-bit truth
table code (outputs for inputs $(0,0), (0,1), (1,0), (1,1)$).

\begin{table}[htbp]
\centering
\caption{Complete table of all 16 two-input Boolean functions.}
\label{tab:all-16-full}
\begin{tabular}{@{}ccccccll@{}}
  \toprule
  \textbf{Index} & $(0,0)$ & $(0,1)$ & $(1,0)$ & $(1,1)$
    & \textbf{Code} & \textbf{Name} & \textbf{Expression} \\
  \midrule
  0  & 0 & 0 & 0 & 0 & 0000 & FALSE     & $0$ \\
  1  & 0 & 0 & 0 & 1 & 0001 & AND       & $A \wedge B$ \\
  2  & 0 & 0 & 1 & 0 & 0010 & Inhibit B & $A \wedge \bar{B}$ \\
  3  & 0 & 0 & 1 & 1 & 0011 & Project A & $A$ \\
  4  & 0 & 1 & 0 & 0 & 0100 & Inhibit A & $\bar{A} \wedge B$ \\
  5  & 0 & 1 & 0 & 1 & 0101 & Project B & $B$ \\
  6  & 0 & 1 & 1 & 0 & 0110 & XOR       & $A \oplus B$ \\
  7  & 0 & 1 & 1 & 1 & 0111 & OR        & $A \vee B$ \\
  8  & 1 & 0 & 0 & 0 & 1000 & NOR       & $\overline{A \vee B}$ \\
  9  & 1 & 0 & 0 & 1 & 1001 & XNOR      & $\overline{A \oplus B}$ \\
  10 & 1 & 0 & 1 & 0 & 1010 & NOT B     & $\bar{B}$ \\
  11 & 1 & 0 & 1 & 1 & 1011 & Imply A   & $A \vee \bar{B}$ \\
  12 & 1 & 1 & 0 & 0 & 1100 & NOT A     & $\bar{A}$ \\
  13 & 1 & 1 & 0 & 1 & 1101 & Imply B   & $\bar{A} \vee B$ \\
  14 & 1 & 1 & 1 & 0 & 1110 & NAND      & $\overline{A \wedge B}$ \\
  15 & 1 & 1 & 1 & 1 & 1111 & TRUE      & $1$ \\
  \bottomrule
\end{tabular}
\end{table}


% ----------------------------------------------------------------------------
\section{Properties and Relationships}
\label{sec:gate-properties}
% ----------------------------------------------------------------------------

\paragraph{Complementary pairs.}
Each function $g$ has a complement $\bar{g} = 15 - g$ (bitwise NOT
of the truth table):
\begin{center}
\begin{tabular}{@{}ll@{}}
  AND (1) $\leftrightarrow$ NAND (14) &
  OR (7) $\leftrightarrow$ NOR (8) \\
  XOR (6) $\leftrightarrow$ XNOR (9) &
  FALSE (0) $\leftrightarrow$ TRUE (15) \\
  Project A (3) $\leftrightarrow$ NOT A (12) &
  Project B (5) $\leftrightarrow$ NOT B (10) \\
  Inhibit B (2) $\leftrightarrow$ Imply B (13) &
  Inhibit A (4) $\leftrightarrow$ Imply A (11) \\
\end{tabular}
\end{center}

\paragraph{Symmetric functions.}
Functions where swapping $A$ and $B$ does not change the output:
AND (1), OR (7), XOR (6), XNOR (9), NAND (14), NOR (8), FALSE (0),
TRUE (15). The remaining 8 functions are asymmetric.

\paragraph{Functionally complete subsets.}
\begin{itemize}
  \item $\{\text{NAND}\}$ alone is functionally complete.
  \item $\{\text{NOR}\}$ alone is functionally complete.
  \item $\{\text{AND}, \text{NOT}\}$ is functionally complete.
  \item $\{\text{AND}, \text{OR}, \text{NOT}\}$ is the classical
    complete set.
\end{itemize}

The LGN~\cite{petersen2022difflogic} uses all 16 functions, which is more than sufficient for
completeness. This redundancy is beneficial: having all gate types
available means the network can represent any function with fewer
layers than a NAND-only circuit.


% ----------------------------------------------------------------------------
\section{Implementation as Lookup Tables}
\label{sec:gate-lut}
% ----------------------------------------------------------------------------

Each gate type is implemented as a 4-entry lookup table stored in
the low 4 bits of a single byte:

\begin{verbatim}
const GATE_LUT: [u8; 16] = [
    0b0000, // 0:  FALSE
    0b0001, // 1:  AND
    0b0010, // 2:  A AND NOT B
    0b0011, // 3:  A
    0b0100, // 4:  NOT A AND B
    0b0101, // 5:  B
    0b0110, // 6:  XOR
    0b0111, // 7:  OR
    0b1000, // 8:  NOR
    0b1001, // 9:  XNOR
    0b1010, // 10: NOT B
    0b1011, // 11: A OR NOT B
    0b1100, // 12: NOT A
    0b1101, // 13: NOT A OR B
    0b1110, // 14: NAND
    0b1111, // 15: TRUE
];

fn eval_gate(gate_type: u8, a: bool, b: bool) -> bool {
    let index = ((a as u8) << 1) | (b as u8);
    (GATE_LUT[gate_type as usize] >> index) & 1 == 1
}
\end{verbatim}

This evaluates any of the 16 gate types in constant time with two
shifts, one OR, one AND, and one comparison---no branching required.
